\documentclass[12pt]{article}
\usepackage[utf8]{inputenc}
\usepackage[margin=1in]{geometry}
\usepackage{amsmath}
\usepackage{graphicx}
\usepackage{booktabs}
\usepackage{hyperref}
\usepackage[round]{natbib}

\title{Potassium Regulates Axon-Oligodendrocyte Signaling and Metabolic Coupling in White Matter}
\author{Anonymous Authors}
\date{}

\begin{document}
\maketitle

\begin{abstract}
Oligodendrocytes provide metabolic support to myelinated axons through the delivery of lactate and pyruvate via monocarboxylate transporters, but how oligodendrocytes detect axonal activity and adjust their metabolic output remains poorly understood. Here we show that potassium released during axonal spiking activates inwardly rectifying Kir4.1 channels on oligodendrocytes, triggering calcium influx that regulates metabolic coupling. Using an ex vivo optic nerve preparation combining electrophysiology with two-photon imaging of genetically encoded sensors, we demonstrate that high-frequency axonal firing (10, 25, 50~Hz) produces frequency-dependent, biphasic calcium responses in mature oligodendrocyte somas ($p < 0.0001$ for 5 vs.\ 30~mM K$^+$). These responses depend on barium-sensitive Kir channels and extracellular calcium influx, with reverse-mode Na$^+$/Ca$^{2+}$ exchanger (NCX) contributing to the response, while AMPA/kainate and purinergic receptor antagonists have no effect. Oligodendrocyte-specific Kir4.1 knockout mice (cKO) show reduced GLUT1 and MCT1 expression in myelin by 2--3 months of age and develop axonal pathology by 7--8 months. Using axonally expressed FRET sensors for lactate and glucose, we find lower resting lactate, reduced activity-induced lactate surges, and drastically reduced glucose consumption rates in cKO axons (control: $2.7 \pm 0.6$\%/min vs.\ cKO: $0.1 \pm 0.2$\%/min, $p = 0.0002$). These findings establish potassium-mediated signaling through Kir4.1 as a critical pathway by which oligodendrocytes sense axonal activity and regulate the metabolic support of myelinated axons.
\end{abstract}

\section{Introduction}

Myelinating oligodendrocytes serve dual functions in the central nervous system: they form the myelin sheath that enables rapid saltatory conduction, and they provide essential metabolic support to the axons they ensheath \citep{nave2010myelination, funfschilling2012glycolytic}. This metabolic support takes the form of glycolytic end-products---primarily lactate and pyruvate---delivered to the periaxonal space through monocarboxylate transporter 1 (MCT1) expressed in the myelin membrane \citep{lee2012oligodendroglia, philips2017oligodendroglia}. Disruption of this metabolic support leads to axonal degeneration even when myelin structure appears intact, establishing oligodendroglial metabolism as a critical determinant of axonal health.

Despite the importance of oligodendrocyte-to-axon metabolic coupling, the signaling mechanisms that allow oligodendrocytes to detect axonal activity and adjust their metabolic output remain poorly understood. Recent work demonstrated that oligodendroglial NMDA receptors respond to axonal glutamate release and regulate glucose transporter (GLUT1) surface expression \citep{saab2016oligodendroglial}. However, glutamatergic signaling represents only one potential communication pathway, and the rapid, activity-dependent adjustment of metabolic support likely involves multiple mechanisms.

Potassium is a promising candidate signal for axon-to-oligodendrocyte communication. During action potential propagation, potassium is released into the extracellular space at nodes of Ranvier and along unmyelinated segments, producing transient elevations in local K$^+$ concentration that scale with firing frequency. Oligodendrocytes express high levels of the inwardly rectifying potassium channel Kir4.1, which is critical for K$^+$ buffering in white matter \citep{neusch2001kir41, larson2018oligodendrocytes}. However, whether Kir4.1-mediated potassium sensing in oligodendrocytes also serves a metabolic signaling function---linking axonal activity to the regulation of metabolic support---has not been tested.

Here we use an ex vivo optic nerve preparation that combines compound action potential (CAP) recording with two-photon imaging of genetically encoded calcium, ATP, lactate, and glucose sensors to comprehensively characterize the K$^+$-oligodendrocyte signaling pathway and its consequences for axonal metabolism \citep{imamura2009visualization, san2013imaging}.

\section{Methods}

\subsection{Mouse Lines}

PLP-CreERT;RCL-GCaMP6s mice provided oligodendrocyte-specific calcium imaging (tamoxifen at 6--8 weeks, experiments 4--12 weeks post-injection). Kir4.1$^{\text{fl/fl}}$;MOGiCre mice (cKO) enabled oligodendrocyte-specific Kir4.1 knockout, compared to Kir4.1$^{\text{fl/fl}}$ littermate controls. Axonal metabolite dynamics were measured using intravitreally delivered AAV expressing FRET sensors: ATeam1.03 for ATP, Laconic for lactate, and FLIIP for glucose (Table~\ref{tab:mice}).

\begin{table}[ht]
\centering
\caption{Mouse lines and genetic tools.}
\label{tab:mice}
\begin{tabular}{lll}
\toprule
Tool & Purpose & Delivery \\
\midrule
PLP-CreERT;RCL-GCaMP6s & OL Ca$^{2+}$ imaging & Tamoxifen \\
Kir4.1$^{\text{fl/fl}}$;MOGiCre & OL-specific Kir4.1 KO & Genetic \\
AAV-ATeam1.03 & Axonal ATP sensor & Intravitreal \\
AAV-Laconic & Axonal lactate sensor & Intravitreal \\
AAV-FLIIP & Axonal glucose sensor & Intravitreal \\
\bottomrule
\end{tabular}
\end{table}

\subsection{Ex Vivo Optic Nerve Preparation}

Optic nerves were dissected and placed in a recording chamber with continuous perfusion. Compound action potentials were evoked by suction electrode stimulation at baseline (0.4~Hz) and during high-frequency trains (10, 25, 50~Hz for 30~seconds). Two-photon imaging was performed simultaneously to capture oligodendrocyte calcium dynamics or axonal metabolite levels.

\subsection{Pharmacology}

The following agents were used: TTX (1~$\mu$M; Na$^+$ channel blocker), BaCl$_2$ (Kir channel blocker), NBQX (50~$\mu$M; AMPA/kainate antagonist), PPADS (100~$\mu$M; P2 receptor antagonist), KB-R7943 (NCX inhibitor), NaN$_3$ (5~mM; mitochondrial complex IV inhibitor), and calcium-free extracellular solution.

\subsection{Myelin Proteomics}

Optic nerves from 2.5-month-old cKO ($n = 4$) and control ($n = 4$) mice were analyzed using TMT-labeled LC-MS/MS \citep{thompson2019tmt} to quantify protein expression changes in myelin fractions.

\section{Results}

\subsection{Frequency-Dependent Calcium Responses in Oligodendrocytes}

High-frequency axonal stimulation produced biphasic calcium responses in oligodendrocyte somas: an initial rise during stimulation followed by a transient undershoot. Response amplitude (area under the curve) scaled with stimulation frequency, with progressively larger responses at 10, 25, and 50~Hz (108 cells from 8 nerves). TTX completely abolished the response, confirming dependence on action potentials. Removal of extracellular calcium strongly diminished the response, indicating a calcium influx mechanism.

\subsection{Potassium as the Signaling Mediator}

Bath application of elevated K$^+$ (5, 10, 30~mM above baseline) for 30 seconds produced dose-dependent calcium elevations in oligodendrocytes (Table~\ref{tab:calcium}), mimicking the stimulus-evoked response. Barium, which blocks Kir channels including Kir4.1, diminished both stimulus-evoked and K$^+$-evoked calcium responses in oligodendrocytes without affecting axonal calcium transients. NBQX ($p = 0.7277$) and PPADS had no significant effect on oligodendrocyte calcium responses, ruling out glutamatergic and purinergic signaling. KB-R7943, an inhibitor of reverse-mode NCX, partially reduced the calcium response, suggesting that NCX contributes to calcium influx following Kir4.1-mediated depolarization.

\begin{table}[ht]
\centering
\caption{K$^+$ dose-dependent Ca$^{2+}$ responses in oligodendrocytes.}
\label{tab:calcium}
\begin{tabular}{lccc}
\toprule
$\Delta$K$^+$ (mM) & Response & Comparison & $p$-value \\
\midrule
5 & Small & 5 vs.\ 10 & 0.0048 \\
10 & Moderate & 5 vs.\ 30 & $< 0.0001$ \\
30 & Large & 10 vs.\ 30 & $< 0.0001$ \\
\bottomrule
\end{tabular}
\end{table}

\subsection{Kir4.1 cKO: Electrophysiology and Pathology}

Kir4.1 cKO mice showed normal baseline CAP peak latencies ($p > 0.76$ for all three peaks) but delayed CAP recovery after high-frequency stimulation, equivalent to the effect of barium application. This confirms impaired K$^+$ clearance in the absence of oligodendroglial Kir4.1 \citep{larson2018oligodendrocytes}. Kir4.1 cKO mice developed axonal pathology by 7--8 months of age, preceded by metabolic changes detectable at 2--3 months.

\subsection{Myelin Proteomics}

TMT-labeled LC-MS/MS of optic nerve myelin from 2.5-month-old mice revealed reduced expression of GLUT1 (Slc2a1) and MCT1 (Slc16a1) in cKO myelin compared to controls (Table~\ref{tab:proteomics}). These changes precede the onset of axonopathy, suggesting that metabolic transporter downregulation is an early consequence of disrupted K$^+$ signaling rather than a secondary effect of axonal damage.

\begin{table}[ht]
\centering
\caption{Key protein expression changes in cKO myelin (2.5 months).}
\label{tab:proteomics}
\begin{tabular}{llll}
\toprule
Protein & Gene & Direction in cKO & Category \\
\midrule
GLUT1 & Slc2a1 & Downregulated & Glucose transporter \\
MCT1 & Slc16a1 & Downregulated & Monocarboxylate transporter \\
\bottomrule
\end{tabular}
\end{table}

\subsection{Axonal Lactate Dynamics}

Using the Laconic FRET sensor expressed in optic nerve axons, cKO mice showed lower resting lactate levels in 10~mM glucose, reduced responses to 20~mM exogenous lactate, and diminished activity-induced lactate surges during 50~Hz stimulation compared to controls. These findings indicate that both basal lactate delivery and activity-dependent lactate supply from oligodendrocytes are impaired when Kir4.1-mediated signaling is disrupted.

\subsection{Axonal Glucose Dynamics}

The FLIIP glucose sensor revealed reduced resting glucose levels and attenuated glucose deprivation responses in cKO axons compared to controls. Activity-induced glucose consumption rates during 50~Hz stimulation were drastically different between genotypes (Table~\ref{tab:glucose}): control axons consumed glucose at $2.7 \pm 0.6$\%/min, while cKO axons showed essentially no activity-induced glucose consumption ($0.1 \pm 0.2$\%/min; $p = 0.0002$). This reveals, for the first time, that oligodendrocytes actively regulate axonal glucose metabolism, and that this regulation depends on Kir4.1.

\begin{table}[ht]
\centering
\caption{Activity-induced axonal glucose consumption during 50~Hz stimulation.}
\label{tab:glucose}
\begin{tabular}{lccc}
\toprule
Measure & Control ($n = 11$) & cKO ($n = 16$) & $p$ \\
\midrule
Consumption rate (\%/min) & $2.7 \pm 0.6$ & $0.1 \pm 0.2$ & 0.0002 \\
\bottomrule
\end{tabular}
\end{table}

\section{Discussion}

We establish potassium as a key signal mediating activity-dependent communication between axons and oligodendrocytes, with Kir4.1 channels serving as the oligodendroglial sensor. The proposed signaling cascade proceeds as follows: action potential propagation releases K$^+$ into the periaxonal space; elevated K$^+$ activates oligodendroglial Kir4.1 channels, producing membrane depolarization; depolarization triggers Ca$^{2+}$ influx partly through reverse-mode NCX; the resulting Ca$^{2+}$ signal adjusts metabolic support by regulating GLUT1 and MCT1 expression, lactate supply, and glucose metabolism \citep{funfschilling2012glycolytic, lee2012oligodendroglia, battefeld2016myelinating}.

The K$^+$-Kir4.1-Ca$^{2+}$ pathway operates in parallel with the previously described glutamate-NMDA receptor pathway \citep{saab2016oligodendroglial}. Our pharmacological results demonstrate that AMPA/kainate and purinergic signaling do not contribute significantly to the acute oligodendrocyte calcium response to axonal activity, at least in the optic nerve preparation. The finding that barium blocks the response while glutamate and purine receptor antagonists do not suggests that K$^+$ is the dominant fast signal mediating activity-dependent coupling in this white matter tract.

The early downregulation of GLUT1 and MCT1 in cKO myelin---before the onset of axonopathy---provides a causal chain linking disrupted K$^+$ signaling to impaired metabolic transporter expression. GLUT1 mediates glucose uptake into oligodendrocytes, while MCT1 mediates lactate export to the periaxonal space. Reduced expression of both transporters would simultaneously compromise glucose supply to the oligodendrocyte and lactate delivery to the axon, consistent with the metabolic phenotype we observe in axonal FRET sensor experiments.

The near-complete loss of activity-induced glucose consumption in cKO axons is a striking finding. It indicates that the metabolic response of axons to their own spiking activity is not cell-autonomous but depends on signaling from ensheathing oligodendrocytes. This positions the axon-oligodendrocyte unit as a metabolic partnership in which the oligodendrocyte functions as a metabolic sensor and regulator, adjusting both its own metabolism and that of the axon it supports \citep{nave2010myelination, philips2017oligodendroglia}.

The late-onset axonopathy (7--8 months) in cKO mice, following months of subclinical metabolic deficiency, resembles the progressive axonal degeneration seen in neurological disorders characterized by oligodendrocyte dysfunction. This temporal pattern suggests that chronic metabolic insufficiency, rather than acute energy failure, drives axonal pathology---a mechanism that could contribute to axonal loss in conditions such as multiple sclerosis, leukodystrophies, and age-related white matter degeneration.

\section{Conclusion}

Potassium released during axonal spiking activates Kir4.1 channels on oligodendrocytes, triggering a calcium signaling cascade that regulates metabolic coupling between glia and axons. Loss of this signaling pathway leads to early downregulation of metabolic transporters, impaired axonal lactate and glucose metabolism, and eventual axonal degeneration. These findings identify the K$^+$-Kir4.1 pathway as a critical link between neural activity and metabolic support in white matter, with implications for understanding axonal vulnerability in demyelinating and neurodegenerative diseases.

\bibliographystyle{plainnat}
\bibliography{references}

\end{document}
